\subsection{Loads and Scope}
The booster is treated as a rigid body. Material failure, structural sizing, and thermal protection are outside this study, so the dynamic pressure, aerodynamic moments, specific force, and heating in Sections~\ref{sec:flight-exposure}--\ref{sec:heating} are requirements that a real vehicle would have to meet, not pass-or-fail results. The touchdown limits concern position, relative velocity, attitude, rotation rate, and propellant reserve. The engine's thrust and gimbal ranges are enforced, but their rates of change are not.

\subsection{Earth Curvature and Rotation}
Over the 278 km range of the flat-Earth design, the flat ground plane departs from the sea surface by about 6.1 km, far more than the 5 m landing tolerance. We therefore solve the spherical equations, Eqs.~\eqref{eq:sphere}--\eqref{eq:spherecontact}, in both the optimization and the replay, rather than correcting the planar landing point afterwards. The two optimal ship positions differ by 5.8 km.

Earth rotation is neglected. Its effect depends on the latitude and heading of the flight, which the problem does not give, and including it would require a three-dimensional trajectory.
